% GENERATED FILE. Edit canonical lessons, then run npm run generate:books.
% canonical-source-hash: fnv1a64:27a3c2184de17dd1
% canonical-lessons: TA-C308-paraparappu, TA-C308-vatanti, TA-C308-kicukicu, TA-C308-ampalappatuttu, TA-C308-pakirankamaka

\chapter{Rumours and a Stir}
\label{ch:ta-rumours-and-a-stir}
\hlchaptermodality{308}

\begin{chapteropening}
\textbf{By the end of this chapter:} \emph{I can talk about rumours and gossip.}

Complete the last lesson of chapter 308: I can talk about rumours and gossip.
\end{chapteropening}

\section[\texorpdfstring{\ta{பரபரப்பு} (paraparappu)}{பரபரப்பு (paraparappu)}]{\ta{பரபரப்பு} (paraparappu) — a sensation; a stir; excitement}

\label{lesson:TA-C308-paraparappu}



\begin{quote}
Before the new one: say the Tamil for a titbit, then the Tamil for to announce.
\end{quote}

\subsection*{You'll want to know: \ta{பரபரப்பு}}
\textbf{\ta{பரபரப்பு}} — \emph{paraparappu} — \textquotedblleft{}a sensation; a stir; excitement\textquotedblright{}.

\begin{quote}
\ta{அந்தச்} \ta{செய்தியால்} \ta{ஊர்} \ta{முழுவதும்} \ta{பரபரப்பு} \ta{இருந்தது}.
\end{quote}

\begin{quote}
\emph{antac ceytiyāl ūr muḻuvatum paraparappu iruntatu.}
\end{quote}

\begin{quote}
Because of that news there was a stir all over town.
\end{quote}

With -\ta{ஆன} it describes a noun, as \ta{ஈரம்} gives \ta{ஈரமான}, damp: \ta{பரபரப்பான} \ta{செய்தி}, sensational news.

\begin{cousinweb}
\textbf{\ta{பரபரப்பு}} (\emph{paraparappu}) looks like an expressive doubled word, of the kind Tamil makes for bustle and hurry; its origin is not settled in the sources at hand.
\end{cousinweb}

\begin{grammarlens}[title={What you've built}]
One more word for telling what happened.
\end{grammarlens}

\subsection*{Your turn}
\begin{itemize}
\raggedright
  \item \emph{Say it:} \emph{paraparappu}
  \item \emph{Say it:} \emph{paraparappu}, once more
  \item \emph{Say it:} \emph{antac ceytiyāl ūr muḻuvatum paraparappu iruntatu.} — Because of that news there was a stir all over town.
  \item \emph{Recall:} say the Tamil for a titbit, then the Tamil for to announce, then say \emph{paraparappu} again
\end{itemize}

\begin{tcolorbox}[breakable,colback=teal!4,colframe=teal!35!black,title={Before you move on}]
What does \ta{பரபரப்பு} mean? (\textquotedblleft{}A sensation; a stir; excitement\textquotedblright{}.) Say it once more.
\end{tcolorbox}
\section[\texorpdfstring{\ta{வதந்தி} (vatanti)}{வதந்தி (vatanti)}]{\ta{வதந்தி} (vatanti) — a rumour}

\label{lesson:TA-C308-vatanti}



\begin{quote}
Before the new one: say the Tamil for to announce, then the Tamil for a sensation.
\end{quote}

\subsection*{You'll want to know: \ta{வதந்தி}}
\textbf{\ta{வதந்தி}} — \emph{vatanti} — \textquotedblleft{}a rumour\textquotedblright{}.

\begin{quote}
\ta{அது} \ta{வெறும்} \ta{வதந்தி}, \ta{உண்மை} \ta{இல்லை}.
\end{quote}

\begin{quote}
\emph{atu veṟum vatanti, uṇmai illai.}
\end{quote}

\begin{quote}
It is only a rumour, not the truth.
\end{quote}

\ta{வதந்தி} \ta{பரவியது}, a rumour spread.

\begin{cousinweb}
\textbf{\ta{வதந்தி}} (\emph{vatanti}) is usually traced to Sanskrit \emph{vadanti}, \textquotedblleft{}they say\textquotedblright{}, a form of \emph{vadati}, \textquotedblleft{}speaks, says\textquotedblright{} (CDIAL 11245).
\end{cousinweb}

\begin{grammarlens}[title={What you've built}]
One more word for telling what happened.
\end{grammarlens}

\subsection*{Your turn}
\begin{itemize}
\raggedright
  \item \emph{Say it:} \emph{vatanti}
  \item \emph{Say it:} \emph{vatanti}, once more
  \item \emph{Say it:} \emph{atu veṟum vatanti, uṇmai illai.} — It is only a rumour, not the truth.
  \item \emph{Recall:} say the Tamil for to announce, then the Tamil for a sensation, then say \emph{vatanti} again
\end{itemize}

\begin{tcolorbox}[breakable,colback=teal!4,colframe=teal!35!black,title={Before you move on}]
What does \ta{வதந்தி} mean? (\textquotedblleft{}A rumour\textquotedblright{}.) Say it once more.
\end{tcolorbox}
\section[\texorpdfstring{\ta{கிசுகிசு} (kicukicu)}{கிசுகிசு (kicukicu)}]{\ta{கிசுகிசு} (kicukicu) — to gossip; to whisper}

\label{lesson:TA-C308-kicukicu}



\begin{quote}
Before the new one: say the Tamil for a sensation, then the Tamil for a rumour.
\end{quote}

\subsection*{You'll want to know: \ta{கிசுகிசு}}
\textbf{\ta{கிசுகிசு}} — \emph{kicukicu} — \textquotedblleft{}to gossip; to whisper\textquotedblright{}.

\begin{quote}
\ta{அவர்கள்} \ta{அவனைப்} \ta{பற்றிக்} \ta{கிசுகிசுத்தார்கள்}.
\end{quote}

\begin{quote}
\emph{avarkaḷ avaṉaip paṟṟik kicukicuttārkaḷ.}
\end{quote}

\begin{quote}
They gossiped about him.
\end{quote}

Past (I …): \ta{கிசுகிசுத்தேன்}, \emph{kicukicuttēṉ}.

As a noun, \ta{கிசுகிசு} is a piece of gossip.

\begin{cousinweb}
\textbf{\ta{கிசுகிசு}} (\emph{kicukicu}) sounds like an imitative doubled word for a whisper; its origin is not settled in the sources at hand.
\end{cousinweb}

\begin{grammarlens}[title={What you've built}]
One more word for telling what happened.
\end{grammarlens}

\subsection*{Your turn}
\begin{itemize}
\raggedright
  \item \emph{Say it:} \emph{kicukicu}
  \item \emph{Say it:} \emph{kicukicu}, once more
  \item \emph{Say it:} \emph{avarkaḷ avaṉaip paṟṟik kicukicuttārkaḷ.} — They gossiped about him.
  \item \emph{Recall:} say the Tamil for a sensation, then the Tamil for a rumour, then say \emph{kicukicu} again
\end{itemize}

\begin{tcolorbox}[breakable,colback=teal!4,colframe=teal!35!black,title={Before you move on}]
What does \ta{கிசுகிசு} mean? (\textquotedblleft{}To gossip; to whisper\textquotedblright{}.) Say it once more.
\end{tcolorbox}
\section[\texorpdfstring{\ta{அம்பலப்படுத்து} (ampalappaṭuttu)}{அம்பலப்படுத்து (ampalappaṭuttu)}]{\ta{அம்பலப்படுத்து} (ampalappaṭuttu) — to expose}

\label{lesson:TA-C308-ampalappatuttu}



\begin{quote}
Before the new one: say the Tamil for a rumour, then the Tamil for to gossip.
\end{quote}

\subsection*{You'll want to know: \ta{அம்பலப்படுத்து}}
\textbf{\ta{அம்பலப்படுத்து}} — \emph{ampalappaṭuttu} — \textquotedblleft{}to expose\textquotedblright{}.

\begin{quote}
\ta{ஒரு} \ta{நிருபர்} \ta{அந்த} \ta{ரகசியத்தை} \ta{அம்பலப்படுத்தினார்}.
\end{quote}

\begin{quote}
\emph{oru nirupar anta rakaciyattai ampalappaṭuttiṉār.}
\end{quote}

\begin{quote}
A reporter exposed that secret.
\end{quote}

Past (I …): \ta{அம்பலப்படுத்தினேன்}, \emph{ampalappaṭuttiṉēṉ}.

\ta{அம்பலம்}, a public place: to bring into the open.

\begin{cousinweb}
\textbf{\ta{அம்பலப்படுத்து}} (\emph{ampalappaṭuttu}) is built on native \ta{அம்பலம்}, \textquotedblleft{}an open space for the public, the village assembly\textquotedblright{} (DEDR 173), with \ta{படுத்து}, \textquotedblleft{}to make\textquotedblright{}: to make public.
\end{cousinweb}

\begin{grammarlens}[title={What you've built}]
One more word for telling what happened.
\end{grammarlens}

\subsection*{Your turn}
\begin{itemize}
\raggedright
  \item \emph{Say it:} \emph{ampalappaṭuttu}
  \item \emph{Say it:} \emph{ampalappaṭuttu}, once more
  \item \emph{Say it:} \emph{oru nirupar anta rakaciyattai ampalappaṭuttiṉār.} — A reporter exposed that secret.
  \item \emph{Recall:} say the Tamil for a rumour, then the Tamil for to gossip, then say \emph{ampalappaṭuttu} again
\end{itemize}

\begin{tcolorbox}[breakable,colback=teal!4,colframe=teal!35!black,title={Before you move on}]
What does \ta{அம்பலப்படுத்து} mean? (\textquotedblleft{}To expose\textquotedblright{}.) Say it once more.
\end{tcolorbox}
\section[\texorpdfstring{\ta{பகிரங்கமாக} (pakiraṅkamāka)}{பகிரங்கமாக (pakiraṅkamāka)}]{\ta{பகிரங்கமாக} (pakiraṅkamāka) — openly; publicly}

\label{lesson:TA-C308-pakirankamaka}



\begin{quote}
Before the new one: say the Tamil for to gossip, then the Tamil for to expose.
\end{quote}

\subsection*{You'll want to know: \ta{பகிரங்கமாக}}
\textbf{\ta{பகிரங்கமாக}} — \emph{pakiraṅkamāka} — \textquotedblleft{}openly; publicly\textquotedblright{}.

\begin{quote}
\ta{அவர்} \ta{தன்} \ta{தவறைப்} \ta{பகிரங்கமாக} \ta{ஒப்புக்கொண்டார்}.
\end{quote}

\begin{quote}
\emph{avar taṉ tavaṟaip pakiraṅkamāka oppukkoṇṭār.}
\end{quote}

\begin{quote}
He openly admitted his mistake.
\end{quote}

-\ta{ஆக} makes the adverb; \ta{பகிரங்கம்}, public.

\begin{cousinweb}
\textbf{\ta{பகிரங்கமாக}} (\emph{pakiraṅkamāka}) is built on \ta{பகிரங்கம்}, \textquotedblleft{}public\textquotedblright{}, borrowed from Sanskrit \emph{bahiraṅga}, \textquotedblleft{}outer, external\textquotedblright{} (Monier-Williams), with the adverb ending -\ta{ஆக}.
\end{cousinweb}

\begin{grammarlens}[title={What you've built}]
One more word for telling what happened.
\end{grammarlens}

\subsection*{Your turn}
\begin{itemize}
\raggedright
  \item \emph{Say it:} \emph{pakiraṅkamāka}
  \item \emph{Say it:} \emph{pakiraṅkamāka}, once more
  \item \emph{Say it:} \emph{avar taṉ tavaṟaip pakiraṅkamāka oppukkoṇṭār.} — He openly admitted his mistake.
  \item \emph{Recall:} say the Tamil for to gossip, then the Tamil for to expose, then say \emph{pakiraṅkamāka} again
\end{itemize}

\begin{tcolorbox}[breakable,colback=teal!4,colframe=teal!35!black,title={Before you move on}]
What does \ta{பகிரங்கமாக} mean? (\textquotedblleft{}Openly; publicly\textquotedblright{}.) Say it once more.
\end{tcolorbox}
